\clearpage
\VerdierPhysicalPageAnchor{0041}%
\VerdierSemanticAnchor{ch1-1-3-6-eq2}
との組で、次の図式が可換となるものである：

\begin{equation}
  \begin{tikzcd}[column sep=7em,row sep=3em]
    FT(hg) \arrow[r,"m(hg)"] \arrow[d,equal] &
      T'(hg)F \arrow[d,equal] \\
    FT(g)T(h) \arrow[d,"m(g)\mathbin{\star}T(h)"'] &
      T'(g)T'(h)F \\
    T'(g)FT(h) \arrow[ur,"T'(g)\mathbin{\star}m(h)"'] &
  \end{tikzcd}
  \qquad g,h\in G.
  \tag{1.3.6.2}
\end{equation}

\VerdierSemanticAnchor{ch1-1-3-6-remark}
$G^\circ$ が圏に狭義に作用していても、関手が $G^\circ$ の作用と狭義に可換
するとは限らないことに注意せよ。したがって $G^\circ$ が圏 $\mathcal C$ と
$\mathcal C'$ に自明に作用する場合（すなわち $T(g)=T'(g)=$ 恒等関手）、
$G^\circ$-関手とは自己同型群における $G$ の表現を備えた関手にほかならない。
また、$G^\circ$-関手が狭義 $G^\circ$-関手と同型になるとは限らない。
[19] に着想を得れば、反例を容易に構成できる。

\VerdierSemanticAnchor{ch1-1-3-6-strictify}
同文献の結果により、任意の $G^\circ$-圏が狭義 $G^\circ$-圏と
$G^\circ$-同値であることも示せる。同じ方法で次のように強められる。
二つの $G^\circ$-圏 $(\mathcal C,G^\circ,T,c)$ と
$(\mathcal C',G^\circ,T',c')$ に対し、二つの狭義 $G^\circ$-圏
$(\widetilde{\mathcal C},G^\circ,\widetilde T)$,
$(\widetilde{\mathcal C}',G^\circ,\widetilde T')$ と $G^\circ$-同値

\VerdierSemanticAnchor{ch1-1-3-6-equivalences}
\[
  (u,\mu)\colon
  (\widetilde{\mathcal C},G^\circ,\widetilde T)
  \longrightarrow
  (\mathcal C,G^\circ,T,c),
\]
\[
  (v,\nu)\colon
  (\mathcal C',G^\circ,T',c')
  \longrightarrow
  (\widetilde{\mathcal C}',G^\circ,\widetilde T')
\]
が存在し、任意の関手

\VerdierSemanticAnchor{ch1-1-3-6-functor}
\[
  (F,m)\colon
  (\mathcal C,G^\circ,T,c)
  \longrightarrow
  (\mathcal C',G^\circ,T',c'),
\]
に対して合成関手

\VerdierSemanticAnchor{ch1-1-3-6-composite}
\[
  (v,\nu)(F,m)(u,\mu)\colon
  (\widetilde{\mathcal C},G^\circ,\widetilde T)
  \longrightarrow
  (\widetilde{\mathcal C}',G^\circ,\widetilde T')
\]
は $G^\circ$ の作用と狭義に可換する関手と同型になる。この結果は用いない。

\VerdierSemanticAnchor{ch1-1-3-7}
\noindent\textbf{1.3.7.}\quad
本小節では、次数付き圏の言語と $G^\circ$-圏の言語という二つの言語の同値性を
振り返った。本書では
